% Artifact Description — IEEE EDGE 2026
% This document follows the IEEE conference formatting guidelines.
% It is used solely for the review process and will NOT appear in proceedings.

\documentclass[conference]{IEEEtran}

\usepackage[utf8]{inputenc}
\usepackage[T1]{fontenc}
\usepackage{hyperref}
\usepackage{booktabs}
\usepackage{listings}
\usepackage{xcolor}
\usepackage{enumitem}

% Code listing style
\lstset{
  basicstyle=\ttfamily\small,
  backgroundcolor=\color{gray!8},
  frame=single,
  framerule=0.4pt,
  rulecolor=\color{gray!40},
  breaklines=true,
  columns=fullflexible,
  keepspaces=true,
  showstringspaces=false,
  xleftmargin=4pt,
  xrightmargin=4pt,
  aboveskip=6pt,
  belowskip=6pt,
}

\begin{document}

\title{Artifact Description:\\KaaS-Edge: Resource-Aware Knowledge Distillation Service for Heterogeneous Wireless Edge Networks}

\author{\IEEEauthorblockN{Anonymous}
\IEEEauthorblockA{Double-blind review}}

\maketitle

% ============================================================
\section{Access \& Links}
% ============================================================

The complete source code, configuration files, and experiment scripts are provided in this artifact package.
The repository is also available at:

\begin{itemize}
  \item \textbf{Repository:} \url{https://anonymous.4open.science/r/Kaas-0DD0}
\end{itemize}

\noindent All materials are released under the MIT License.

% ============================================================
\section{User Guide}
% ============================================================

\subsection{Installation}

\begin{lstlisting}[language=bash]
git clone <REPO_URL>
cd KaaS-Edge
python -m venv venv
source venv/bin/activate   # Linux / macOS
pip install -r requirements.txt
\end{lstlisting}

\subsection{Reproducing Main Results (Table~I \& Figures~2--3)}

Run the main comparison experiment across three random seeds:

\begin{lstlisting}[language=bash]
python scripts/run_edge_experiments.py \
    --exp main --device cuda:0
\end{lstlisting}

Or use parallel execution for faster runs:

\begin{lstlisting}[language=bash]
bash parallel_run.sh main cuda:0
python merge_seed_results.py
\end{lstlisting}

\subsection{Reproducing Sensitivity Analysis (Figures~4--5)}

\textbf{Budget sensitivity (Figure~4):}

\begin{lstlisting}[language=bash]
python scripts/run_edge_experiments.py \
    --exp budget --device cuda:0
\end{lstlisting}

\textbf{Device scalability:}

\begin{lstlisting}[language=bash]
python scripts/run_edge_experiments.py \
    --exp scale --device cuda:0
\end{lstlisting}

\textbf{Privacy impact (Figure~5):}

\begin{lstlisting}[language=bash]
python scripts/run_edge_experiments.py \
    --exp privacy --device cuda:0
\end{lstlisting}

\subsection{Running All Experiments at Once}

\begin{lstlisting}[language=bash]
python scripts/run_edge_experiments.py \
    --all --device cuda:0
\end{lstlisting}

\subsection{Generating Figures}

After experiments complete:

\begin{lstlisting}[language=bash]
python plot_edge_figures.py
\end{lstlisting}

Figures are saved to the \texttt{figures/} directory.

\subsection{Quick Verification (CPU-Only)}

For reviewers without GPU access, a quick smoke-test validates basic executability:

\begin{lstlisting}[language=bash]
python scripts/run_edge_experiments.py \
    --exp main --quick
\end{lstlisting}

This runs with reduced sample sizes on CPU and completes in approximately 5--10~minutes.

% ============================================================
\section{Environment}
% ============================================================

\subsection{Software Dependencies}

Table~\ref{tab:deps} lists the key software dependencies.
A complete list is provided in \texttt{requirements.txt}.

\begin{table}[h]
\centering
\caption{Key Software Dependencies}
\label{tab:deps}
\begin{tabular}{@{}ll@{}}
\toprule
Package & Version \\
\midrule
Python       & $\geq$ 3.8   \\
PyTorch      & $\geq$ 2.0.0 \\
torchvision  & $\geq$ 0.15.0 \\
numpy        & $\geq$ 1.24.0 \\
scipy        & $\geq$ 1.10.0 \\
scikit-learn & $\geq$ 1.3.0 \\
matplotlib   & $\geq$ 3.7.0 \\
pyyaml       & $\geq$ 6.0   \\
tqdm         & $\geq$ 4.65.0 \\
\bottomrule
\end{tabular}
\end{table}

\subsection{Hardware Used for Paper Results}

\begin{table}[h]
\centering
\caption{Hardware Configuration}
\label{tab:hw}
\begin{tabular}{@{}ll@{}}
\toprule
Component & Specification \\
\midrule
GPU     & NVIDIA RTX 5070 Ti (16\,GB VRAM) \\
CPU     & Intel i7 equivalent \\
RAM     & $\geq$ 16\,GB \\
Storage & $\geq$ 5\,GB free \\
\bottomrule
\end{tabular}
\end{table}

\subsection{Minimum Requirements}

\begin{table}[h]
\centering
\caption{Minimum Hardware Requirements}
\label{tab:minreq}
\begin{tabular}{@{}ll@{}}
\toprule
Task & Minimum Hardware \\
\midrule
Quick test (\texttt{--quick}) & CPU only, 8\,GB RAM \\
Single experiment             & GPU with $\geq$ 4\,GB VRAM \\
Full reproduction             & GPU with $\geq$ 8\,GB VRAM \\
\bottomrule
\end{tabular}
\end{table}

% ============================================================
\section{Data \& Methods}
% ============================================================

\subsection{Dataset}

\textbf{CIFAR-100} (Krizhevsky, 2009) is automatically downloaded by \texttt{torchvision} on first run.
It contains 50{,}000 training images and 10{,}000 test images across 100 classes (32$\times$32 pixels).
The training set is split into private data (40k) and public reference data (10k) using a fixed random seed, ensuring zero data leakage.

\subsection{Data Partitioning}

Non-IID data distribution across devices is created using Dirichlet partitioning with concentration parameter $\alpha = 0.5$ (configurable in \texttt{config/edge\_default.yaml}).

\subsection{Device Simulation}

Heterogeneous edge devices are simulated with log-normal cost distributions.
Each device profile includes:
\begin{itemize}[noitemsep]
  \item \textbf{Cost} ($b_i$): marginal per-sample upload cost
  \item \textbf{Privacy} ($\rho_i$): privacy degradation factor $\in (0, 1]$
  \item \textbf{Capacity} ($\theta_i$): half-saturation constant
\end{itemize}
Device generation is deterministic given a random seed, ensuring full reproducibility.

% ============================================================
\section{Assessment: Interpreting Results}
% ============================================================

\subsection{Output Format}

Each experiment produces a JSON file in \texttt{results/edge/} containing per-round metrics including accuracy history, final accuracy, total cost, and number of rounds.

\subsection{Key Metrics}

\begin{table}[h]
\centering
\caption{Key Evaluation Metrics}
\label{tab:metrics}
\begin{tabular}{@{}lll@{}}
\toprule
Metric & Description & Expected Range \\
\midrule
\texttt{final\_accuracy} & Test accuracy after 50 rounds & 44--51\% \\
\texttt{total\_cost}     & Cumulative resource cost      & $\leq B \times T$ \\
\texttt{n\_participants} & Devices selected per round     & 5--15 \\
\texttt{comm\_mb}        & Per-round communication (MB)   & 7--15 \\
\bottomrule
\end{tabular}
\end{table}

\subsection{Expected Outcomes}

The experiments should demonstrate:

\begin{enumerate}[noitemsep]
  \item \textbf{Main comparison (Table~I, Figure~2):} KaaS-Edge achieves accuracy comparable to full-participation baselines while reducing per-round communication by $\sim$9.7$\times$ and cumulative bandwidth by $\sim$17$\times$.
  \item \textbf{Budget sensitivity (Figure~4):} Accuracy increases with budget $B$ following a concave curve, confirming the submodular structure of the quality function.
  \item \textbf{Device scalability:} Quality improves with more devices due to the submodular greedy selection exploiting device diversity.
  \item \textbf{Privacy impact (Figure~5):} Higher privacy (lower $\rho$) degrades accuracy gradually; KaaS-Edge maintains 41.3\% under strong privacy ($\bar{\rho} \approx 0.1$), remaining competitive with baselines at weaker privacy.
\end{enumerate}

\subsection{Variance}

All experiments use 3 random seeds (42, 123, 456).
Results are reported as mean $\pm$ standard deviation.
Minor numerical differences across hardware are expected but should not affect relative rankings between methods.

% ============================================================
\section{System Requirements}
% ============================================================

\subsection{Full Reproduction}

\begin{itemize}[noitemsep]
  \item \textbf{GPU:} NVIDIA GPU with $\geq$ 8\,GB VRAM and CUDA support
  \item \textbf{Time:} $\sim$2--3 hours for all experiments on a modern GPU
  \item \textbf{Disk:} $\sim$2\,GB for CIFAR-100 download + results
\end{itemize}

\subsection{Simplified Verification}

For reviewers with limited resources:

\begin{lstlisting}[language=bash]
python scripts/run_edge_experiments.py \
    --exp main --quick
\end{lstlisting}

This validates that: (i) all dependencies install correctly, (ii) the RADS scheduler produces valid allocations, (iii) the training loop completes without errors, and (iv) results are saved in the expected format.

% ============================================================
\section{File Inventory}
% ============================================================

Table~\ref{tab:files} provides a summary of the key files in the repository.

\begin{table}[h]
\centering
\caption{Key Repository Files}
\label{tab:files}
\begin{tabular}{@{}p{0.42\columnwidth}p{0.50\columnwidth}@{}}
\toprule
File / Directory & Description \\
\midrule
\texttt{src/scheduler/rads.py}      & Core RADS algorithm (Proposition~1 + Theorem~1) \\
\texttt{src/methods/kaas\_edge.py}  & KaaS-Edge method implementation \\
\texttt{src/methods/fed*.py}        & Baseline method implementations \\
\texttt{src/privacy/ldp.py}         & Local differential privacy mechanisms \\
\texttt{scripts/run\_edge\_experiments.py} & Main experiment runner \\
\texttt{config/edge\_default.yaml}  & Default configuration \\
\texttt{parallel\_run.sh}           & Multi-seed parallel execution \\
\texttt{merge\_seed\_results.py}    & Result aggregation across seeds \\
\texttt{plot\_edge\_figures.py}     & Figure generation script \\
\texttt{requirements.txt}           & Python package dependencies \\
\bottomrule
\end{tabular}
\end{table}

\end{document}
